\documentclass[10pt,leqno,twoside]{article}

% Custom journal/conference class (defines \fofej, page layout, etc.)
\usepackage{annal-latex}

% Math
\usepackage{amsmath}

% Graphics & figures
\usepackage{graphicx}
\usepackage{tikz}
\usepackage{booktabs}  % \toprule / \midrule / \bottomrule used in the Results table

% Watermark & color (for draft marking)
\usepackage{draftwatermark}
\usepackage{color}

% Bibliography
\usepackage[backend=bibtex, sortcites=true, style=numeric-comp]{biblatex}
\addbibresource{second.bib}

% Author formatting
\usepackage{authblk}

\SetWatermarkText{\textbf{FAKE RESEARCH}}
\SetWatermarkLightness{0.85}
\SetWatermarkAngle{60}
\SetWatermarkScale{3.0}

% Reusable author/affiliation block for the closing page, to avoid
% repeating the same 6 lines four times with copy-paste risk.
\newcommand{\authorblock}[2]{%
  \noindent\textbf{#1}\\
  Department of Computer Science\\Eötvös Loránd University\\
  Pázmány Péter Sétány 1/C Budapest, Hungary\\
  Budapest\\
  Hungary\\
  {\tt #2}\\
  \\
}

\begin{document}
\setcounter{page}{1}

% \fofej{volume}{issue}{start-page-code}{end-page-code}
\fofej{49}{19}{nn}{nnn}

\vspace{.4cm}

% ---------------------------------------------------------------

\title{The Not For You Page:\\Algorithmic Design for User Disengagement}

\author{%
  {\bf Fodor Máté} (Budapest, Hungary)\\[1ex]
  {\bf Merk Eszter} (Budapest, Hungary)\\
  {\bf Vidó Bence} (Budapest, Hungary)\\
  {\bf Preszter Balázs} (Budapest, Hungary)%
}

\date{September 2026}

% ---------------------------------------------------------------


\abstract{Algorithmic recommendation systems across modern digital platforms are predominantly engineered to maximize user engagement, retention, and time on platform. While effective for commercial metrics, these continuous feed architectures frequently lead to digital fatigue, compulsive consumption, and diminished user agency. This paper explores the concept of the ``Not For You Page'' (NFYP), an alternative feed model designed to investigate how recommendation systems can be re-oriented toward disengagement and user well-being. By examining mechanisms that introduce intentional friction and alter content distribution, this research evaluates how algorithmic design can encourage voluntary platform exit. The project addresses the feasibility of embedding digital well-being directly into recommendation logic, offering a counter perspective to traditional retention focused platform architectures.}

% ---------------------------------------------------------------

\section{Introduction}
Feed-based recommendation engines serve as the foundational interface for modern content discovery. Driven by retention metrics, these architectures continuously adapt to user behavior to capture and sustain attention. However, the systematic optimization of screen time increasingly conflicts with user autonomy, giving rise to digital burnout and cognitive overload. Current interventions for digital well-being rely almost exclusively on external controls, such as application timers or system level access restrictions, leaving the underlying algorithms largely unexamined as tools for self-regulation.

To address this issue, this project examines the ``Not For You Page'' (NFYP), a conceptual recommendation model that inverts conventional retention objectives. Rather than reinforcing engagement loops, the system explores methods to systematically reduce content resonance and introduce friction as usage continues. This research investigates the theoretical foundation of anti-retention algorithm design, outlining potential mechanisms for disengagement and discussing their broader implications for human-computer interaction and ethical platform development.

% ---------------------------------------------------------------

\subsection{Related work}

Research on TikTok's recommendation system has documented how its algorithm is explicitly engineered around engagement maximization. Zhou~\cite{zhou2024understanding} shows that TikTok blends collaborative and content-based filtering to personalize content delivery, but finds that this comes at the cost of content homogeneity and low algorithmic transparency.

This engagement-first design has been linked to a range of documented harms. Smith and Short~\cite{smith2022needs} provide psychometric evidence of problematic TikTok use, with withdrawal and relapse emerging as the strongest predictive criteria and loneliness positively associated with maladaptive use patterns. Beyond behavioral addiction, Ionescu and Licu~\cite{ionescu2023tiktok} review evidence that, over time, algorithmically curated feeds shape users' self-perceived identity and personal values. At a societal level, Kuncoro and Hasanah~\cite{kuncoro2024social} argue that the same engagement-maximizing logic can create echo chambers that reinforce radical or extremist views by continuously amplifying content that sustains user attention regardless of its ideological content.

In parallel, human-computer interaction research has explored interventions for digital well-being, but almost exclusively at the interface layer rather than within the recommendation logic itself. Haliburton et al.~\cite{haliburton2024friction} show that short, interface-level design frictions can reduce compulsive app-opening behavior over time, while Almoallim and Sas~\cite{almoallim2022wellbeing} find that the most popular digital well-being applications rely almost entirely on screen-time monitoring and externally imposed usage limits, without modifying the underlying recommendation algorithm.

Taken together, this literature establishes both that TikTok-style algorithms are explicitly optimized for engagement and that this optimization carries documented costs to individual behavior, identity, and public discourse. Yet the interventions proposed to date treat well-being as something to protect through external, interface-level controls, never by reorienting the recommendation objective itself. The Not For You Page addresses this gap directly.

% ---------------------------------------------------------------

\section{Methodology}

\subsection{Research Design and System Mechanics}
This study evaluates an anti-recommender system designated as the Not For You Page, engineered to minimize screen time and prompt voluntary user disengagement. The primary research question addresses how algorithmic filtering can reliably shorten user sessions without employing obvious attrition measures. Instead, the system targets standard social media content that is inherently underwhelming, uneventful, or lacking resolution, regardless of its production value or format.

The recommendation engine parses a user's historical interaction data to identify primary interest vectors and systematically suppresses those topics. Concurrently, the algorithm samples various categories of unengaging media, evaluating how effectively each post accelerates session termination. Through continuous real-time feedback, the system refines its selection strategy to isolate which specific patterns of uneventful content trigger the fastest exits.

\subsection{Participants and Experimental Setup}
The system was evaluated with a sample of forty volunteers recruited through internal announcements. To qualify, participants were required to be at least eighteen years of age and active daily users of social media platforms. All individuals were informed that they were evaluating a novel feed interface, but they remained naive to the disengagement hypothesis to ensure their exit behavior reflected natural usage patterns.

Testing was conducted in a quiet environment using a custom web application designed to replicate a standard mobile social media client. The software automatically logged interaction telemetry, including scroll velocity, item view duration, and precise timestamps for session initiation and termination. Participants accessed the feed on standardized mobile devices and were instructed to browse naturally until they chose to stop.

\subsection{Procedure and Data Collection}
Upon arrival, participants completed an initial profiling assessment to establish baseline content preferences across major media categories. They were then randomly assigned in a 1:1 ratio to either the control condition or the treatment condition. The control group received a standard recommendation feed optimized for content relevance and retention, while the treatment group received the Not For You Page feed.

Data collection commenced upon the initial swipe and concluded when the participant explicitly closed the application interface. The primary outcome variable was total session duration, measured in seconds from feed entry to exit. Secondary metrics included average dwell time per post, total items viewed, and the frequency of rapid-skip interactions. Following session termination, participants completed a brief post-test questionnaire assessing their subjective motivation for exiting the application.

\subsection{Data Analysis Plan}
The primary analytical objective was to evaluate whether the Not For You Page produced a statistically significant reduction in total session duration relative to the control feed. Independent samples t-tests were used to compare mean session lengths and per-post view times between the two groups. Additionally, survival analysis using Kaplan-Meier estimates was applied to compare retention probability curves over time and evaluate overall disengagement velocity. Secondary interaction data was analyzed to determine which specific categories of underwhelming content generated the highest skip rates prior to session abandonment.

% ---------------------------------------------------------------


\section{Results}
\begin{table}[ht]
\centering
\begin{tabular}{lrr}
\toprule
A & B & C \\
\midrule
A1 & B1 & C1\\
A2 & B2 & C2\\
A3 & B3 & C3\\
A4 & B4 & C4\\
\bottomrule
\end{tabular}
\caption{\label{tab:train}The summary of the table}
\end{table}

Reference to the Table \ref{tab:train} on page \pageref{tab:train} and a cite \cite{reversing}.

% \begin{figure}
%     \centering
%     \includegraphics[width=\linewidth]{dog}
%     \caption{Cute golden retriever puppies}
%     \label{fig:dog}
% \end{figure}

% ---------------------------------------------------------------


\section{Discussion} % Conclusions, Further research
Your discussion

% ---------------------------------------------------------------

\section*{Acknowledgment}
Your acknowledgement

% ---------------------------------------------------------------


\printbibliography[title=References]


\clearpage
\vspace{2cm}

\authorblock{Fodor Máté}{email}
\authorblock{Merk Eszter}{merkeszter@gmail.com}
\authorblock{Vidó Bence}{email}
\authorblock{Preszter Balázs}{email}

\end{document}